\begin{figure}[!htbp]
  \centering
  \begin{tikzpicture}[
    every node/.style={font=\small, align=center},
    box/.style={rectangle, rounded corners=3pt, draw=blue!55!black, thick,
                fill=blue!8, minimum width=2.2cm, minimum height=0.9cm},
    warn/.style={rectangle, rounded corners=3pt, draw=red!65!black, thick,
                 fill=red!8, minimum width=2.4cm, minimum height=0.9cm},
    ok/.style={rectangle, rounded corners=3pt, draw=green!45!black, thick,
               fill=green!10, minimum width=2.2cm, minimum height=0.9cm},
    side/.style={rectangle, rounded corners=3pt, draw=gray!70, thick,
                 fill=gray!10, minimum width=2.5cm, minimum height=1.3cm},
    arr/.style={-stealth, thick}
  ]
    \node[font=\bfseries] at (0,3.1) {naive dequant + attention};
    \node[box]  (na1) at (-4.6,2.0) {低比特 KV cache};
    \node[warn] (na2) at (-1.8,2.0) {物化完整 FP16 K/V};
    \node[warn] (na3) at (1.0,2.0) {物化完整 attention matrix};
    \node[box]  (na4) at (3.8,2.0) {标准 attention};
    \draw[arr] (na1.east) -- (na2.west);
    \draw[arr] (na2.east) -- (na3.west);
    \draw[arr] (na3.east) -- (na4.west);

    \node[font=\bfseries] at (0,0.7) {Triton fused decode};
    \node[ok] (fu1) at (-5.0,-0.3) {低比特 K/V + metadata};
    \node[ok] (fu2) at (-2.5,-0.3) {block load};
    \node[ok] (fu3) at (0.0,-0.3) {SRAM / register\\unpack + dequant};
    \node[ok] (fu4) at (2.6,-0.3) {$qK^\top$ + online softmax};
    \node[ok] (fu5) at (5.2,-0.3) {accumulate output};
    \draw[arr] (fu1.east) -- (fu2.west);
    \draw[arr] (fu2.east) -- (fu3.west);
    \draw[arr] (fu3.east) -- (fu4.west);
    \draw[arr] (fu4.east) -- (fu5.west);

    \node[side] (gqa) at (6.0,1.55) {\textbf{GQA mapping}\\
      $N_{\mathrm{rep}}=H_q/H_{kv}$\\
      $h_{kv}=\lfloor h_q/N_{\mathrm{rep}}\rfloor$};
    \draw[arr] (gqa.south) |- (fu4.north);
  \end{tikzpicture}
  \caption{Triton 融合解码核的数据流：低比特缓存、片上反量化与在线 softmax。图中对比 naive 路径与 fused 路径的关键差异：后者避免物化完整 FP16 K/V 与完整 attention matrix，并将 GQA 头映射直接纳入 decode 数据流。}
  \label{fig:ch3-triton-dataflow}
\end{figure}
